\subsection{Extension of representations to spaces of measures}

In this section, we assume that K = C.

Let \(\varrho\) be a continuous linear representation of G in a separated locally convex quasi-complete space E. Let us denote by \(\mathcal{M}_{c}(\mathrm{G})\) the space of measures with compact support on G equipped with the topology of compact convergence; it is the dual of the space \(\mathcal{C}(\mathrm{G})\) . For every measure \(\nu \in \mathcal{M}_{c}(\mathrm{G})\) we set

\[
\varrho (\nu) = \int_ {\mathrm{G}} \varrho (g) d \nu (g).
\]

This is an element of \(\mathcal{L}(\mathrm{E})\) . In particular, we have \(\varrho(\varepsilon_g) = \varrho(g)\) for all \(g \in \mathrm{G}\) .

According to INT, VIII, p. 136, § 2, no. 6, the map \(\nu \mapsto \varrho(\nu)\) is a continuous linear map from \(\mathcal{M}_c(\mathrm{G})\) to the space \(\mathcal{L}(\mathrm{E})\) equipped with the topology of compact convergence. According to INT, VIII, p. 145, § 3, no. 3, prop. 11, it is a morphism of unital algebras.

Let \(x \in \mathrm{E}\) . The map from \(\mathcal{M}_c(\mathrm{G})\) to E defined by \(\nu \mapsto \varrho(\nu)x\) is continuous when \(\mathcal{M}_c(\mathrm{G})\) is equipped with the topology of compact convergence (INT, VI, p. 27, § 1, no. 7, prop. 16 and EVT, III, p. 31, prop. 4).

For \(f \in \mathcal{L}^{1}(\mathrm{G})\) vanishing outside a compact subset of G, the measure \(f \cdot \mu\) has compact support and we shall denote by \(\varrho^{\mu}(f)\) , or \(\varrho(f)\) when no ambiguity is possible, the endomorphism \(\varrho(f \cdot \mu)\) of E. This endomorphism depends only on the class \(\widetilde{f}\) of \(f\) in \(\mathrm{L}^{1}(\mathrm{G})\) , and we shall also denote it by \(\varrho^{\mu}(\widetilde{f})\) or \(\varrho(\widetilde{f})\) .

Similarly, let \(\varrho\) be a continuous bounded linear representation of G in a Banach space E. For every bounded measure \(\nu \in \mathcal{M}^{1}(\mathrm{G})\) we set

\[
\varrho (\nu) = \int_ {\mathrm{G}} \varrho (g) d \nu (g).
\]

According to INT, VIII, loc. cit., the map \(\nu \mapsto \varrho(\nu)\) is a continuous unital morphism of Banach algebras from the algebra \(\mathcal{M}^{1}(G)\) of bounded complex measures on G into the Banach algebra \(\mathcal{L}(E)\) .

Let \(\rho\) be the function \(g\mapsto\|\varrho(g)\|\) on G; the algebra denoted \(\mathcal{M}^{\rho}\) in INT, VIII, p. 145, prop. 11 (whose elements are the measures \(\nu\) such that \(\rho\in\mathcal{L}^{1}(\nu)\) ) coincides with the Banach algebra \(\mathcal{M}^{1}(G)\) . Indeed, set \(M = \sup \rho\) . We have \(M > 0\) since \(\rho(e)=1\) , and \(M^{-1} \leqslant \rho \leqslant M\) since \(\|\varrho(e)\|\leqslant\|\varrho(g)\|\|\varrho(g^{-1})\|\) for all \(g\in G\) ; thus \(\rho\in\mathcal{L}^{1}(\nu)\) if and only if \(\nu\) is a bounded measure.

If \(f \in \mathcal{L}^{1}(\mathrm{G})\) , one also denotes \(\varrho^{\mu}(f)\) or simply \(\varrho(f)\) the endomorphism \(\varrho(f \cdot \mu)\) , and likewise for the class \(\tilde{f}\) of \(f\) in \(\mathrm{L}^{1}(\mathrm{G})\) .

Lemma 1. --- Let \(\varrho\) be a unitary representation of G in a Hilbert space E. The map \(\nu \mapsto \varrho(\nu)\) of \(\mathcal{M}^{1}(\mathrm{G})\) into \(\mathcal{L}(\mathrm{E})\) is a unital morphism of involutive Banach algebras.

By what precedes, it suffices to show that the morphism \(\nu \mapsto \varrho(\nu)\) is involutive. Let \(\nu \in \mathcal{M}^{1}(\mathrm{G})\) . The measure \(\nu^{*}\) is the conjugate measure of the measure \(\check{\nu}\) (I, p. 99, example 4). According to the definition of the conjugate measure (INT, III, p. 52, § 1, no. 5), one computes

\[
\begin{array}{c} \langle x \mid \varrho (\nu) y \rangle = \int_ {\mathrm{G}} \langle x \mid \varrho (g) y \rangle d \nu (g) \\ = \int_ {\mathrm{G}} \langle \varrho (g ^ {- 1}) x \mid y \rangle d \nu (g) = \int_ {\mathrm{G}} \langle \varrho (g) x \mid y \rangle d \check {\nu} (g) \\ = \overline {{\int_ {\mathrm{G}} \langle y \mid \varrho (g) x \rangle d \nu^ {*} (g)}} = \langle \varrho (\nu^ {*}) x \mid y \rangle \end{array}
\]

for all \(x\) and \(y\) in E, whence \(\varrho(\nu)^{*} = \varrho(\nu^{*})\) .

Let \(\varrho\) be a continuous linear representation (resp. continuous and bounded) of G in a locally convex quasi-complete space E (resp. a Banach space E). If F is a closed subspace of E defining a subrepresentation \(\varrho_{\mathrm{F}}\) of \(\varrho\) , then for every measure \(\nu \in \mathcal{M}_c(\mathrm{G})\) (resp. \(\nu \in \mathcal{M}^1 (\mathrm{G})\) ) the linear map \(\varrho (\nu)\) leaves the subspace F stable and coincides with \(\varrho_{\mathrm{F}}(\nu)\) on F.

Conversely, let \(F \subset E\) be a closed subspace, stable under the linear maps \(\varrho(f)\) for every function \(f \in \mathcal{H}(G)\) . Then F defines a subrepresentation of \(\varrho\) (INT, VIII, p. 139, § 2, no. 7, prop. 10).

Recall (A, VIII, p. 388) that a function \(f \colon \mathbf{G} \to \mathbf{C}\) is said to be central if, for all \(g\) and \(h\) in \(\mathbf{G}\) , we have \(f(gh) = f(hg)\) . This amounts to saying that \(f\) is invariant under conjugation.

DEFINITION 1. --- A bounded measure \(\nu \in \mathcal{M}^{1}(\mathrm{G})\) is said to be central if one has \(\varepsilon_{g} * \nu = \nu * \varepsilon_{g}\) for all \(g \in \mathrm{G}\) .

If G is unimodular and \(f \in \mathcal{C}(\mathrm{G})\) is \(\mu\) -integrable, the measure \(f \cdot \mu\) is central if and only if f is central.

Let \(\varrho\) be a continuous bounded representation of G in a Banach space E (resp. a continuous representation of G in a quasi-complete separated locally convex space E). For every bounded central measure \(\nu\) on G (resp. any compactly supported central measure \(\nu\) on G), the linear map \(\varrho(\nu)\) is a G-morphism from \(\varrho\) to \(\varrho\) . Indeed, for every \(g \in G\) , we have

\[
\varrho (\nu) \varrho (g) = \varrho (\nu * \varepsilon_ {g}) = \varrho (\varepsilon_ {g} * \nu) = \varrho (g) \varrho (\nu).
\]

